The CBV core calculus is translated into Call-by-Push-Value (CBPV)
\parencite{levy_2003} before execution. This intermediate representation is not
merely a compilation target: it is the calculus in which the equational theory of
\parencite{jones_2026} is formulated, and whose equations justify both the
transitions of the abstract machine and the rewrite rules of the peephole
optimiser. In this section we describe the translation, present the equational
theory, and show how it validates program transformations.

\subsection{Translation to CBPV}
\label{subsection:translation}

\begin{figure*}
  \begin{footnotesize}
\begin{minipage}{0.28\textwidth}
{\color{CBVGreen}\fbox{\color{black}\text{type translation}}}
\; $\Transl{A}$
\begin{align*}
    \Transl{\UnitT} &= \UnitT \\
    \Transl{\NatT} &= \NatT \\
    \Transl{A_1 + A_2} &= \Transl{A_1} + \Transl{A_2} \\
    \Transl{A_1 \times A_2} &= \Transl{A_1} \times \Transl{A_2} \\
    \Transl{\ListT{A}} &= \ListT{\Transl{A}} \\
    \Transl{A_1 \to A_2} &= \ThunkT*{\Transl{A_1} \to \ReturnerT*{\Transl{A_2}}}
\end{align*}
\end{minipage}
\begin{minipage}{0.70\textwidth}
{\color{CBVGreen}\fbox{\color{black}\text{term translation}}}
\; $\Transl{\CBV{e}} = \CTerm{M}$
\begin{align*}
    \Transl{\CBV{x}} &= \CTerm{\Return{x}}
    \\
    \Transl{\CBV{\UnitV}} &= \CTerm{\Return{\UnitV}}
    \qquad
    \Transl{\CBV{\Zero}} = \CTerm{\Return{\Zero}}
    \\
    \Transl{\CBV{\Succ{e}}} &= \CTerm{\CTo{\Transl*{\CBV{e}}}{v}{\Return{\Succ{v}}}}
    \\
    \Transl{\CBV{\VTuple{e_1}{e_2}}} &= \CTerm{\CTo{\Transl*{\CBV{e_1}}}{v_1}{\CTo*{\Transl*{\CBV{e_2}}}{v_2}{\Return{\VTuple{v_1}{v_2}}}}}
    \\
    \Transl{\CBV{\Nil}} &= \CTerm{\Return{\Nil}}
    \\
    \Transl{\CBV{\Cons{e_1}{e_2}}} &= \CTerm{\CTo{\Transl*{\CBV{e_1}}}{v}{\CTo*{\Transl*{\CBV{e_2}}}{w}{\Return{\Cons{v}{w}}}}}
    \\
    \Transl{\CBV{\CBVLam{x}{e}}} &= \CTerm{\Return*{\Thunk{\Lam{x}{\Transl{\CBV{e}}}}}}
    \\
    \Transl{\CBV{\CBVApp{e_1}{e_2}}} &= \CTerm{\CTo{\Transl*{\CBV{e_1}}}{f}{\CTo*{\Transl*{\CBV{e_2}}}{x}{\CApp{\Force{f}}{x}}}}
    \\
    \Transl{\CBV{\CBVLet{x}{e_1}{e_2}}} &= \CTerm{\CTo{\Transl*{\CBV{e_1}}}{x}{\Transl*{\CBV{e_2}}}}
    \\
    \Transl{\CBV{\CBVIfz{e}{\ldots}{\ldots}{\ldots}}} &= \CTerm{\CTo{\Transl*{\CBV{e}}}{v}{\CIfz{v}{\Transl*{\CBV{e_0}}}{n}{\Transl*{\CBV{e_1}}}}}
    \\
    \Transl{\CBV{\Fail}} &= \CTerm{\Fail}
    \qquad
    \Transl{\CBV{\Choice{}{e_1}{e_2}}} = \CTerm{\Choice{}{\Transl*{\CBV{e_1}}}{\Transl*{\CBV{e_2}}}}
    \\
    \Transl{\CBV{\Exists{x}[\sigma]{e}}} &= \CTerm{\Exists{x}[\sigma]{\Transl*{\CBV{e}}}}
    \\
    \Transl{\CBV{\CBVEquate{e_1}{e_2}{e}}} &= \CTerm{\CTo{\Transl*{\CBV{e_1}}}{v}{\CTo*{\Transl*{\CBV{e_2}}}{w}{\Equate{v}{w}{\Transl*{\CBV{e}}}}}}
\end{align*}
\end{minipage}
\end{footnotesize}

  \caption{Translation from {\color{CBVGreen}CBV} to
    {\color{ACMMatch1}CBPV} {\color{ACMMatch2}terms}. Each CBV expression
    $\CBV{e}$ of type $A$ translates to a CBPV computation $\CTerm{M}$ of
    type $\ReturnerT*{\Transl{A}}$. Constructor arguments are sequenced via
    $\textsf{to}$, making the evaluation order explicit.}
  \label{figure:translation}
\end{figure*}

The translation (\cref{figure:translation}) follows the standard embedding of
call-by-value into CBPV \parencite[\S 4]{levy_2003}, extended with the FLP
effects. Each CBV expression $\CBV{e}$ of type $A$ becomes a CBPV computation of
type $\ReturnerT*{\Transl{A}}$.

The key type translation is for functions: a CBV type $A_1 \to A_2$ becomes the
CBPV value type $\ThunkT*{\Transl{A_1} \to \ReturnerT*{\Transl{A_2}}}$. The
outer $\ThunkT{(-)}$ makes the function a storable value; the inner
$\ReturnerT{(-)}$ marks the result as an effectful computation. All other types
translate homomorphically.

The term translation makes evaluation order explicit. A variable $\CBV{x}$
becomes $\CTerm{\Return{x}}$: a trivially completed computation. A constructor
with subexpressions, such as $\CBV{\Succ{e}}$, must first evaluate its argument:
$\CTerm{\CTo{\Transl*{\CBV{e}}}{v}{\Return{\Succ{v}}}}$. Application
$\CBV{\CBVApp{e_1}{e_2}}$ evaluates both subexpressions, forces the resulting
thunk, and applies it. A $\lambda$-abstraction $\CBV{\CBVLam{x}{e}}$ becomes
$\CTerm{\Return*{\Thunk{\Lam{x}{\Transl{\CBV{e}}}}}}$. The FLP effects
translate directly: $\CBV{\Fail}$ becomes $\CTerm{\Fail}$,
$\CBV{\Choice{}{e_1}{e_2}}$ becomes $\CTerm{\Choice{}{M_1}{M_2}}$, and so on.

\paragraph{Why CBPV?}
The value/computation distinction is essential for three reasons. First, it makes
the sequencing of effects explicit: every point at which an effect may occur is
visible in the CBPV term, which is critical for reasoning about nondeterminism.
Second, CBPV provides the framework in which the equational theory of
\parencite{jones_2026} is formulated; without it, the equations that justify our
optimiser would not be available. Third, the CBPV structure maps directly onto
the abstract machine: values are data, computations are processes, and the
machine's stack tracks pending continuations.

\subsection{The Equational Theory}
\label{subsection:equations}

\begin{figure*}
  \begin{footnotesize}
\fbox{\textbf{$\beta$-laws}} \\[0.5ex]
\begin{minipage}{0.48\textwidth}
\begin{align*}
  \CTerm{\CTo*{\Return{V}}{x}{M}}
    &= \CTerm{\CSubst{M}{V}{x}} \\
  \CTerm{\CApp*{\Lam{x}{M}}{V}}
    &= \CTerm{\CSubst{M}{V}{x}} \\
  \CTerm{\Force*{\Thunk{M}}}
    &= \CTerm{M}
\end{align*}
\end{minipage}
\begin{minipage}{0.48\textwidth}
\begin{align*}
  \CTerm{\CIfz{\Zero}{M}{n}{N}}
    &= \CTerm{M} \\
  \CTerm{\CIfz{\Succ{V}}{M}{n}{N}}
    &= \CTerm{\CSubst{N}{V}{n}} \\
  \CTerm{\CMatch{\Cons{V}{W}}{M}{x}{xs}{N}}
    &= \CTerm{\SubstTwo{N}{V}{W}{x}{xs}}
\end{align*}
\end{minipage}

\bigskip

\fbox{\textbf{$\eta$-laws}} \\[0.5ex]
\begin{minipage}{0.48\textwidth}
\begin{align*}
  \CTerm{M}
    &= \CTerm{\CTo{M}{x}{\Return{x}}} \\
  \CTerm{M}
    &= \CTerm{\Lam{x}{\CApp{M}{x}}}
\end{align*}
\end{minipage}
\begin{minipage}{0.48\textwidth}
\begin{align*}
  \VTerm{V}
    &= \VTerm{\Thunk*{\Force{V}}} \\
  \CTerm{\CSubst{M}{V}{n}}
    &= \CTerm{\CIfz{V}{\CSubst{M}{\Zero}{n}}{m}{\CSubst{M}{\Succ{m}}{n}}}
\end{align*}
\end{minipage}

\bigskip

\fbox{\textbf{Sequencing laws}} \\[0.5ex]
\begin{align*}
  \CTerm{\CTo*{\CTo{M}{x}{N}}{y}{P}}
    &= \CTerm{\CTo{M}{x}{(\CTo{N}{y}{P})}}
  & \CTerm{\CApp*{\CTo{M}{x}{N}}{V}}
    &= \CTerm{\CTo{M}{x}{\CApp{N}{V}}}
\end{align*}

\bigskip

\fbox{\textbf{Effect equations}} \\[0.5ex]
\begin{minipage}{0.48\textwidth}
\begin{align*}
  \CTerm{\Choice{\Fail}{M}} &= \CTerm{M} \\
  \CTerm{\Choice{M}{\Fail}} &= \CTerm{M} \\
  \CTerm{\Choice{M}{N}} &= \CTerm{\Choice{N}{M}} \\
  \CTerm{\Choice{M}{M}} &= \CTerm{M}
\end{align*}
\end{minipage}
\begin{minipage}{0.48\textwidth}
\begin{align*}
  \CTerm{\Equate{V}{V}{M}} &= \CTerm{M} \\
  \CTerm{\CTo{M}{x}{\Fail}} &= \CTerm{\Fail} \\
  \CTerm{\Exists{x}[\sigma]{\Fail}} &= \CTerm{\Fail} \\
  \CTerm{\Choice{M_1}{\DelimPrn{\Choice{M_2}{M_3}}}}
    &= \CTerm{\Choice{\DelimPrn{\Choice{M_1}{M_2}}}{M_3}}
\end{align*}
\end{minipage}

\bigskip

\fbox{\textbf{Pull-through laws}} \\[0.5ex]
\begin{align*}
  \CTerm{\CTo*{\Choice{M_1}{M_2}}{x}{N}}
    &= \CTerm{\Choice{(\CTo{M_1}{x}{N})}{(\CTo{M_2}{x}{N})}}
  \\
  \CTerm{\CTo*{\Exists{z}[\sigma]{M}}{x}{N}}
    &= \CTerm{\Exists*{z}[\sigma]{\CTo{M}{x}{N}}}
  \\
  \CTerm{\CTo*{\Equate{V}{W}{M}}{x}{N}}
    &= \CTerm{\Equate{V}{W}{\CTo{M}{x}{N}}}
\end{align*}

\bigskip

\fbox{\textbf{$\lambda$-laws}} \\[0.5ex]
\begin{align*}
  \CTerm{\Lam{x}{\Fail}}
    &= \CTerm{\Fail}
  &
  \CTerm{\Lam*{x}{\Choice{M}{N}}}
    &= \CTerm{\Choice{\Lam{x}{M}}{\Lam{x}{N}}}
  \\
  \CTerm{\Lam*{x}{\Exists{z}[\sigma]{M}}}
    &= \CTerm{\Exists*{z}[\sigma]{\Lam{x}{M}}}
  &
  \CTerm{\Lam*{x}{\Equate{V}{W}{M}}}
    &= \CTerm{\Equate{V}{W}{\Lam{x}{M}}}
\end{align*}

\fbox{\textbf{Parameter equations} (example: $\NatT$)} \\[0.5ex]
\begin{align*}
  \CTerm{\Exists{n}[\NatT]{M}}
    &= \CTerm{\Choice{\CSubst{M}{\Zero}{n}}{\Exists{m}[\NatT]{\CSubst{M}{\Succ{m}}{n}}}}
\end{align*}

\fbox{\textbf{Narrowing equations} (examples)} \\[0.5ex]
\begin{align*}
  \CTerm{\Equate{\Succ{V}}{\Succ{W}}{M}}
    &= \CTerm{\Equate{V}{W}{M}}
  &
  \CTerm{\Equate{\Succ{V}}{\Zero}{M}}
    &= \CTerm{\Fail}
\end{align*}
\end{footnotesize}

  \caption{Selected equations of the CBPV equational theory, validated by the
    denotational semantics of \parencite{jones_2026}. All equations are between
    well-typed CBPV terms.}
  \label{figure:equations}
\end{figure*}

The denotational semantics of \parencite{jones_2026} validates a rich equational
theory for CBPV with FLP effects. These equations serve two purposes: they
justify the transitions of the abstract machine (\cref{section:machine}), and
they provide the rewrite rules for the peephole optimiser
(\cref{section:optimisation}). We collect the most important equations in
\cref{figure:equations} and discuss them briefly.

\paragraph{Structural equations}
The $\beta$-laws express the computational content of each eliminator: case
analysis on a known constructor reduces, application of a $\lambda$-abstraction
substitutes, forcing a thunk unwraps it, and the bind-return law
$\CTerm{\CTo{\Return{V}}{x}{M}} = \CTerm{\CSubst{M}{V}{x}}$ expresses that
returning a value and immediately binding it is the same as substitution. The
$\eta$-laws express the universality of each type's eliminator. The sequencing
laws express associativity of bind and its interaction with application.

\paragraph{Effect equations}
The algebraic structure of the FLP effects gives rise to a family of equations.
Failure is the identity for choice: $\CTerm{\Choice{\Fail}{M}} = \CTerm{M}$.
Choice is associative, commutative, and idempotent:
$\CTerm{\Choice{M}{M}} = \CTerm{M}$. Existential quantifiers commute with each
other and distribute over choice. Equational constraints are reflexive
($\CTerm{\Equate{V}{V}{M}} = \CTerm{M}$), commute with each other, and
distribute over choice. The \emph{dead-end} law states that
$\CTerm{\CTo{M}{x}{\Fail}} = \CTerm{\Fail}$: if the continuation always fails,
the preceding computation is irrelevant.

\paragraph{Pull-through laws}
The \emph{pull-through} laws describe how effects interact with sequencing. If
the bound computation branches, the continuation runs in each branch:
$\CTerm{\CTo*{\Choice{M_1}{M_2}}{x}{N}} =
\CTerm{\Choice{(\CTo{M_1}{x}{N})}{(\CTo{M_2}{x}{N})}}$. Existential
quantifiers and equational constraints can be pulled through a bind:
$\CTerm{\CTo*{\Exists{z}[\sigma]{M}}{x}{N}} =
\CTerm{\Exists*{z}[\sigma]{\CTo{M}{x}{N}}}$. These laws are used extensively by
the optimiser to simplify machine terms.

\paragraph{Parameter and narrowing equations}
The \emph{parameter equations} characterise the behaviour of $\exists$ on each
first-order type. For example, $\CTerm{\Exists{n}[\NatT]{M}} =
\CTerm{\Choice{\CSubst{M}{\Zero}{n}}{\Exists{m}[\NatT]{\CSubst{M}{\Succ{m}}{n}}}}$:
introducing a logical variable of type $\NatT$ is the same as choosing between
zero and the successor of a fresh variable. The \emph{narrowing equations} govern
unification: matching constructors are peeled off
($\CTerm{\Equate{\Succ{V}}{\Succ{W}}{M}} = \CTerm{\Equate{V}{W}{M}}$) and
mismatched constructors cause failure
($\CTerm{\Equate{\Succ{V}}{\Zero}{M}} = \CTerm{\Fail}$).

\subsection{CBV Equations}
\label{subsection:cbv-equations}

\begin{figure}
\begin{small}
\begin{align}
  \CBV{\CBVLet{x}{(\CBVLet{y}{e_1}{e_2})}{e_3}}
    &= \CBV{\CBVLet{y}{e_1}{(\CBVLet{x}{e_2}{e_3})}}
  \label{eq:cbv-assoc} \tag{Assoc} \\
  \CBV{\CBVLet{x}{e}{x}}
    &= \CBV{e}
  \label{eq:cbv-eta} \tag{$\eta$-Let} \\
  \CBV{\CBVLet{x}{\Fail}{e}}
    &= \CBV{\Fail}
  \label{eq:cbv-fail} \tag{Dead-End} \\
  \CBV{\CBVLet{x}{(\Choice{}{e_1}{e_2})}{e}}
    &= \CBV{\Choice{}{(\CBVLet{x}{e_1}{e})}{(\CBVLet{x}{e_2}{e})}}
  \label{eq:cbv-choice} \tag{$\talloblong$-Let} \\
  \CBV{\Choice{}{e}{\Fail}}
    &= \CBV{e}
  \label{eq:cbv-fail-id} \tag{$\talloblong$-Id} \\
  \CBV{\Choice{}{e}{e}}
    &= \CBV{e}
  \label{eq:cbv-idem} \tag{$\talloblong$-Idem}
\end{align}
\end{small}
\caption{Selected CBV equations, valid by translation to CBPV.}
\label{figure:cbv-equations}
\end{figure}

The CBPV equations induce equations on the CBV calculus. If $\Transl{\CBV{e_1}}$
and $\Transl{\CBV{e_2}}$ are equal as CBPV terms (using the equations of
\cref{figure:equations}), then $\CBV{e_1}$ and $\CBV{e_2}$ are
observationally equivalent. \Cref{figure:cbv-equations} lists several such
equations. These are the equations that a programmer can rely on when reasoning
about \Gweek{} programs.

\paragraph{Example proof}
We demonstrate the method by proving \eqref{eq:cbv-assoc}: the associativity of
let-bindings. Translating both sides to CBPV, the left-hand side becomes
$\CTerm{\CTo{(\CTo{\Transl*{\CBV{e_1}}}{y}{\Transl*{\CBV{e_2}}})}{x}{\Transl*{\CBV{e_3}}}}$
and the right-hand side becomes
$\CTerm{\CTo{\Transl*{\CBV{e_1}}}{y}{(\CTo{\Transl*{\CBV{e_2}}}{x}{\Transl*{\CBV{e_3}}})}}$.
These are equal by the CBPV sequencing law
$\CTerm{\CTo*{\CTo{M}{x}{N}}{y}{P}} =
\CTerm{\CTo{M}{x}{(\CTo{N}{y}{P})}}$ from \cref{figure:equations}.

\paragraph{A plausible but invalid equation}
Not every equation that appears plausible is valid. Consider:
%
\[
  \CBV{\CBVLet{x}{e_1}{\Choice{}{e_2}{e_3}}}
  \stackrel{?}{=}
  \CBV{\Choice{}{(\CBVLet{x}{e_1}{e_2})}{(\CBVLet{x}{e_1}{e_3})}}
\]
%
This `right-distribution' of let over choice would state that binding over a
choice is the same as choosing between bindings. In a deterministic language this
would be harmless, but in FLP it is \emph{not} valid: if $\CBV{e_1}$ is
nondeterministic (e.g.\ $\CBV{e_1} = \CBV{\Choice{}{\Zero}{\Succ{\Zero}}}$),
the left-hand side evaluates $\CBV{e_1}$ once and branches on the continuation,
while the right-hand side evaluates $\CBV{e_1}$ twice, potentially doubling the
search tree. The CBPV translation makes this clear: the left-hand side is
$\CTerm{\CTo{M}{x}{\Choice{N_1}{N_2}}}$, which has a single copy of $\CTerm{M}$
above the choice, while the right-hand side is
$\CTerm{\Choice{(\CTo{M}{x}{N_1})}{(\CTo{M}{x}{N_2})}}$, which duplicates
$\CTerm{M}$. These are not equal: the commutativity law (which allows reordering
of independent binds) does not apply because $\CBV{e_2}$ and $\CBV{e_3}$ may
depend on $x$. The idempotence of choice
($\CTerm{\Choice{M}{M}} = \CTerm{M}$) does not help either, because the
duplicated copies may produce different \emph{numbers} of results, not just
different results.
